% !TeX program = xelatex
% !BIB program = biber

\documentclass[a4paper, monochrome]{article}

\usepackage{longtable}
\usepackage{array}
\usepackage{placeins}
\usepackage{csquotes}
\usepackage{hyperref}
\usepackage{graphicx}
\usepackage{url}
\usepackage{todonotes}
\usepackage{booktabs}
\usepackage{tablefootnote}
\usepackage{dirtree}
\usepackage{siunitx}
\usepackage{multirow}
\usepackage{tabularx}
\sisetup{output-exponent-marker=\ensuremath{\mathrm{e}}}
\usepackage[ngerman]{babel}
\usepackage[
    backend=biber,
    style=apa,
    autocite=inline,
    mincitenames=1,
    maxcitenames=2
]{biblatex}

\addbibresource{thesis-literature.bib}
\addbibresource{thesis-software.bib}
\addbibresource{thesis-online.bib}

\newcommand{\code}[1]{\texttt{#1}}

\selectlanguage{ngerman}

\def\UrlFont{\rmfamily}

\begin{document}

\begin{titlepage}
    \begin{center}
        \vspace*{1cm}

        \textbf{\Large Bildbasierte Digitalisierung von historischen Toninformationsträgern}

        \vspace{0.5cm}
        Untertitel
                
        \vspace{1.5cm}


        \vfill
                
        Master-Arbeit zur Erlangung des akademischen Grades eines\\
        M.Sc. Digital Humanities

        \vspace{3cm}
                
        Eingereicht von: David Fuhry
        
        \vspace{0.2cm}

        Matrikel-Nr. 3704472\\
        Zschochersche Str. 32\\
        04229 Leipzig

        \vspace{1cm}

        Erstbegutachter: Prof. Dr. Phil. Manuel Burghardt \\
        Zweitbegutachter: N.N.
        
        \vspace{0.2cm}
                    
        Computational Humanities\\
        Institut für Informatik\\
        Universität Leipzig\\
        Augustusplatz 10\\
        04109 Leipzig 
                
    \end{center}
\end{titlepage}

\tableofcontents

\thispagestyle{empty}

\newpage

% Introduction

\pagenumbering{arabic}

\section{Einleitung}

Historische Medien wie Lochplatten und Klavierrollen sind die ältesten, kommerziell in größerer Stückzahl produzierten und auch für den Privatgebrauch gedachten Toninformationsträger.
Sie wurden vom Ende des 19. Jahrhunderts bis zur frühen Mitte des 20. Jahrhunderts von einer Vielzahl von Musikstücken, durch verschiedene Hersteller und in einer hohen Zahl unterschiedlicher Formate produziert.
Die Medien sind zeithistorische Zeugnisse aus einer Zeit, in der privater Musikkonsum erstmals einer breiteren Gruppe von Personen möglich wurde.

Neben durch die Hersteller aus Kompositionen erstellten Medien existieren mit den sogenannte Player Rolls auch Notenrollen, die an einem speziell modifizierten Klavier von kontemporären Pianist:innen eingespielt wurden.
Diese Notenrollen bilden die einzigen Aufnahmen vieler der bekanntesten Pianist:innen aus einer Zeit, in der Audioaufnahmen technisch noch nicht ausgereift genug für die Aufnahme musikalischer Darbietungen waren.

Diese Eigenschaften machen solche historischen Toninformationsträger zu einer wichtigen Ressource für die moderne musikwissenschaftliche Forschung.
Während eine signifikante Anzahl von Medien, insbesondere Notenrollen, bis heute überlebt hat, wird sich ihr Zustand auf absehbare Zeit weiter verschlechtern und so zu irreversiblen Verlusten von Medien und den darauf enthaltenen Informationen führen.

Um diese zeithistorischen Dokumente auch in Zukunft für die musikwissenschaftliche Forschung zugänglich zu und gleichzeitig weitere Schäden an den Medien, etwa durch Abnutzung, zu minimieren bietet sich die Digitalisierung ebendieser an.
Dazu zählt zum einen die Erstellung von hochauflösenden Bildern der Medien, zum anderen aber auch die Digitalisierung der auf den Medien kodierten Toninformationen.
Verfahren zu diesem Zweck existieren, sind aber überwiegend schlecht dokumentiert und der Öffentlichkeit nicht problemlos zugänglich (siehe Abschnitt \ref{Forschungsstand}).

Das Musikinstrumentenmuseum der Universität Leipzig ist im Besitz einer umfangreichen Sammlung von solchen historischen Toninformationsträgern.
Im Rahmen des vergangenen Forschungsprojektes TASTEN wurden von ca. 3000 der Notenrollen aus dieser Sammlung Scans erstellt.

Ziel des aktuellen Forschungsvorhabens DISKOS an der Forschungsstelle des Musikinstrumentenmuseums ist die Digitalisierung dieser Notenrollen und weiterer, plattenförmiger Toninformationsträger aus derselben Sammlung, sowie die inhaltliche Analyse der digitalisierten Medien.
Eine kurze Übersicht dazu findet sich bei \textcite[]{khulusi_2022}.
Zur Umsetzung der Ziele des Forschungsprojektes soll unter anderem eine Anwendung entwickelt werden, die Medien aller Formate in der Sammlung verarbeiten kann.
Die Anwendung soll auch für die inhaltlichen Analysen nutzbar sein.

Zu diesem Zweck wurde die in dieser Arbeit beschriebene Software implementiert.
Sie ist in der Lage, aus den Bilddaten von historischen Toninformationsträgern, die darauf kodierten musikalischen Informationen zu extrahieren.

Es wird zunächst ein voll funktionsfähiges Verfahren zur Digitalisierung von Lochplatten vorgestellt.
Ebenso wird eine erste Implementierung für ein Verfahren zur Digitalisierung von Notenrollen eingeführt.
Die Anwendung gibt nach der Verarbeitung der Bilddaten eine standardisierte Midi-Datei aus.
Im weiteren werden die in die Anwendung integrierten Analyse- und Visualisierungstools, die in Zusammenarbeit mit den am Forschungsprojekt beteiligten Musikwissenschaftler:innen zur Beantwortung musikwissenschaftlicher Fragestellungen entwickelt wurden, demonstriert.



% Historische Toninformationsträger

\FloatBarrier
\section{Historische Toninformationsträger}

\begin{figure*}[t]
    \centering
    \includegraphics[width=\textwidth]{graphics/placeholder.png}
    \caption{Lochplatten zweier verschiedener Typen aus dem Bestand des Musikinstrumentenmuseums der Universität Leipzig}
    \label{platten}
\end{figure*}

Als Toninformationsträger\footnote{Auch Historical Music Storage Media.} im Sinne dieser Arbeit werden die Medien bezeichnet, die zum Betrieb von mechanischen Musikinstrumenten benötigt werden.
Eine Definition für mechanische Musikinstrumente ist etwa in \textcite[]{mgg_mechanische} gegeben:

\begin{quotation}
    Mechanische Musikinstrumente (selbstspielende Musikinstrumente, Musikautomaten, Automatophone) sind Musikinstrumente, bei denen die Tonfolge selbsttätig durch einen Toninformationsträger gesteuert wird. Die Mehrzahl der mechanischen Musikinstrumente reproduziert Musik ohne Einwirkungsmöglichkeit des Menschen. \parencite[I. Definition]{mgg_mechanische}
\end{quotation}

Die Geschichte solcher mechanischen Musikinstrumente reicht zurück bis in die Antike, in der bereits einfache Orgeln existierten die über das setzen von Stiften auf einer sich drehenden Walze gesteuert werden konnten.
Den Höhepunkt ihrer Popularität erreichten sie während des späten 19. und frühen 20. Jahrhunderts, als eine Vielzahl von verschiedenen, immer komplexeren mechanischen Instrumenten konstruiert wurde \parencite[10-12]{bowers_1972}.
Mit der Entwicklung von besseren und günstigeren Methoden zur Audioaufnahme und -reproduktion in den 1930ern wurden die mechanischen Musikinstrumente praktisch komplett vom Markt verdrängt \parencite[2]{zoltan_1994}.
Für eine ausführliche Übersicht über die Geschichte der mechanischen Musikinstrumente sei auf die entsprechende Literatur verwiesen, etwa \textcite[]{bowers_1972,bowers_1975} sowie \textcite[]{mgg_mechanische}.

Die Medien für solche mechanischen Instrumente kodieren Informationen aus denen das mechanische Instrument Musik erzeugen kann.
Üblicherweise werden dafür zumindest grundlegende Informationen wie der Beginn, die Länge sowie die zugehörige Note für jeden Ton auf dem Medium kodiert.
Es existieren aber auch Instrumente und mit ihnen Formate von Toninformationsträgern die zusätzlich weitere Informationen zu Dynamik und Spielweise kodieren können.
Im weiteren sollen vorrangig die für diese Arbeit relevanten Typen und Formate von Toninformationsträgern weiter betrachtet werden.
% Platten: 1882 bis 1910, Rollen 1900 bis 1930
Diese wurden im Zeitraum zwischen 1882 und 1930 produziert und lassen sich grob in zwei Kategorien unterteilen: Plattenförmige sowie rollenförmige Medien, wobei die plattenförmigen Medien sowohl zeitlich als auch technologisch als Vorgänger der Rollen betrachtet werden können.

\subsection{Plattenförmige Toninformationsträger}

Plattenförmigen Medien wurden ab ca. 1875 produziert und stellten einen paradigmenwechsel in der Konstruktion mechanischer Musikinstrumente dar.
Waren bei früheren Modellen die Toninformationen direkt im Instrument kodiert (beispielsweise über Walzen), wurden sie nun separat auf Platten kodiert und waren somit erstmals unabhängig vom Instrument.
Eine solche Platte hat zumeist eine kurze Laufzeit von einer kleinen einstelligen Zahl von Minuten und konnte damit zumindest kürzere Stücke oder Motive aus größeren Kompositionen enthalten.
Durch diese Entwicklung wurde es erstmals Möglich eine größere Sammlung von Musikstücken zu besitzen \parencite[III.5.c. Plattenspieldosen und Drehinstrumente]{mgg_mechanische}.

Frühe Formate dieser Form von mechanischen Instrumente nutzten Platten aus Pappe auf denen die Toninformationen in Form von Löchern kodiert sind (im Folgenden Lochplatten genannt).
Lochplatten dieser Formate sind sich grundsätzlich recht ähnlich und unterscheiden sich vor allem in den verfügbaren Tönen sowie den physischen Abmessungen der Medien.
Beispiele für zwei Lochplatten unterschiedlicher Formate dieses Typs sind in Abbildung \ref{platten} zu sehen.

Die Lochplatten haben mehrere Spuren, jede Spur ist einem bestimmten Ton zugeordnet, meist sind die Spuren aufsteigend nach Tonhöhe sortiert, mit der tiefsten Note auf der innersten Spur.
Die Töne selbst wurden mittels Löchern binär in die Platten kodiert, dort wo ein Loch vorhanden ist klingt der Ton der entsprechenden Spur.
Wichtig anzumerken ist, dass die meisten Formate nicht chromatisch angelegt waren, die populären Platten der Marke Ariston beispielsweise haben 24 Spuren decken aber über 3 Oktaven ab.\footnote{Siehe \textcite[]{mxp_2003520} für eine genaue Aufschlüsselung des Tonvorrats dieses Formates.}
Durch das auslassen bestimmter Töne ergaben sich zwar Einschränkungen bei der Kodierung von musikalischen Stücken auf die Medien, es konnte aber ein größerer Tonumfang abgebildet werden ohne die Medien und Abspielgeräte in unpraktikable Größen zu skalieren.

\begin{figure*}[t]
    \centering
    \includegraphics[width=\textwidth]{graphics/placeholder.png}
    \caption{Abspielgerät für Lochplatten der Marke Ariston aus dem Bestand des Musikinstrumentenmuseums der Universität Leipzig}
    \label{aristonplayer}
\end{figure*}

Die Abspielgeräte für solche Platten (siehe Abbildung \ref{aristonplayer}) nutzen Verfahren zur Tonerzeugung, dass dem einer Mundharmonika nicht unähnlich ist.
Eine Lochplatte wurde in ein Abspielgerät eingelegt das über eine Leiste mit Zungen verfügt die den gesamten Radius des informationstragenden Teils der Lochplatte bedeckt.
Zum Abspielen wurde die Lochplatte, je nach Gerät per Handkurbel oder automatisiert, Gedreht und Luft wurde durch die Zungen bewegt.
An den Stellen an denen Löcher vorhanden waren konnte die Luft durch die Zungen passieren, so dass ein Ton erzeugt wurde.
Der so entstehende Klang dieser sogenannten Organetten wurde jedoch häufig als nicht zufriedenstellend empfunden \parencite[III.5.c. Plattenspieldosen und Drehinstrumente]{mgg_mechanische}.

Eine der Entwicklungen die gemacht wurden um dieses Defizit im Klang auszugleichen war der mechanische Klaviervorsetzer.
Dieser nutzt Platten die denen für die Organetten sehr ähnlich sind und Informationen nach dem selben Prinzip kodierten.
Im Gegensatz zu diesen war im Abspielgerät selbst (siehe Abbildung \ref{aristonplayer}) allerdings keine Tonerzeugung integriert.
Stattdessen besaß das Gerät eine Reihe mechanischer Finger, die, vor einem Klavier platziert, das auf der Lochplatte kodierte Stück auf diesem "spielten".

Eine andere Weiterentwicklung der frühen Organetten stellen plattenförmige Medien aus Metall dar.
Bei der Produktion solcher wurden kleine Teile des Blechs umgebogen, so dass Noppen entstanden die in der Lage waren Tonkämme anzureißen.
Die Abspielgeräte für solche Platten verfügten häufig über Zusatzinstrumente wie Triangeln oder Trommeln und waren unter anderem in Wirtshäusern sehr populär \parencite[III.5.c. Plattenspieldosen und Drehinstrumente]{mgg_mechanische}.
Während die Digitalisierung solcher Metallplatten Teil des DISKOS Projektes ist, ist sie aufgrund der fundamental unterschiedlichen Art der Informationskodierung auf diesen Medien kein Teil dieser Arbeit.

Im Rahmen des DISKOS-Projektes wurden Bilder eines vollständiges Konvolut aus 60 Lochplatten der Marke Ariston angefertigt die mit der in dieser Arbeit vorgestellten Software digitalisiert werden.
Weitere 40 Lochplatten für einen Klaviervorsetzer befinden sich noch im Fotografieprozess, konnten aber bereits für Testdigitalisierungen genutzt werden.

\subsection{Notenrollen}

\begin{figure*}[t]
    \centering
    \includegraphics[width=\textwidth]{graphics/placeholder.png}
    \caption{Abspielgerät für Lochplatten der Marke Ariston aus dem Bestand des Musikinstrumentenmuseums der Universität Leipzig}
    \label{pianoroll}
\end{figure*}

Während die mit der Lochplatte eingeführte Trennung von Abspielgerät und Toninformationen eine signifikante Weiterentwicklung gegenüber den früheren, walzenbasierten mechanischen Musikinstrumenten darstellte, wurde mit der Notenrolle\footnote{Auch Klavierrolle genannt.} ein weiterer Entwicklungsschritt gemacht.
Dieses um die Jahrhundertwende eingeführte Format unterscheidet sich zum einen in der Form der Toninformationsträger, die nun als mehrere Meter lange Rollen ausgeführt wurden, aber vor allem durch die eingesetzten Techniken in den Abspielgeräten.
Ein Abbild einer solchen Rolle ist Abbildung \ref{pianoroll} zu entnehmen.

Die Veränderungen an der Art der Informationskodierung auf den Notenrollen im Vergleich zu Lochplatten ist relativ klein.
Zunächst unterscheiden sich die Notenrollen vor allem durch die neue physische Form die eine Spielzeit von bis zu 15 Minuten ermöglichte.
Eine Notenrolle ist analog zu den früheren Lochplatten in mehrere Spuren unterteilt, die nun allerdings nicht mehr kreisförmig, sondern parallel, meist über die Rollenbreite von Bass (Links) nach Diskant (Rechts) sortiert, angelegt sind.
Die Spuren enthalten weiterhin Löcher die über ihre Position die Toninformationen auf der Notenrolle kodieren.
Ein wesentlicher Unterschied zu Pappplatten besteht in der Kodierung länger zu spielender Töne: Diese werden nun nicht mehr durch eine durchgehende Lochung sondern durch in sehr kurzem Abstand platzierte, z.T. auch überlappende, Löcher kodiert.

Frühe Abspielgeräte funktionierten, analog zu mechanischen Klaviervorsetzern für Pappplatten, als Vorsetzer und unterschieden sich im wesentlichen durch die längere Laufzeit der Notenrollen.
Der entscheidende Vorteil dieses Systems lag in der pneumatischen Steuerung des Abspielgeräts.
Zum Abspielen wird die Rolle in das Abspielgerät eingelegt und dieses erzeugte, bei den ersten Geräten durch menschliche Kraft, später durch Elektromotoren, ein Vakuum über die gesamte Rollenbreite.
Erreichte die Rolle nun eine Stelle an der ein Loch vorhanden war, kann die Luft durch das Loch strömen und so eine pneumatische Steuerung betreiben.
Diese erlaubte es über die Stärke des angelegten Vakuums auch die Stärke des Anschlags auf dem bespielten Klavier zu variieren \parencite[III.5.d. Selbstspielende Klaviere]{mgg_mechanische}.

Mit dieser Technik wurden kurz vor der Jahrhundertwende die ersten selbstspielenden Klaviere\footnote{Auch Player Pianos genannt.} konstruiert.
Die meisten dieser Klaviere verfügten über 2 separate Windladen für die Erzeugung des Vakuums und konnten so die Lautstärke von Bass und Diskant getrennt regulieren.
Während die Variation der Vakuumstärke bei den früheren Modellen, genau wie die Wahl der Geschwindigkeit, noch durch die benutzende Person erfolgte, wurden schnell zunehmend komplexere Abspielgeräte und Rollenformate entwickelt bei denen diese Aufgabe vom Abspielgerät selbst übernommen wurde \parencite[III.5.d. Selbstspielende Klaviere]{mgg_mechanische}.

Im Jahr 1904 patentierte die Firma Welte mit dem Welte-Mignon-System das erste Reproduktionsklavier\footnote{Auch Reproducing Player Piano genannt.} ein System, das komplett ohne menschliche Eingaben Musikstücke wiedergeben konnte \parencite[]{mgg_mechanische}.
Um dies zu ermöglichen wurden Informationen beispielsweise zu Dynamik und Pedalbenutzung auf dafür vorgesehenen Spuren der Notenrolle kodiert.
Die Rollen des Typs Welte-Mignon Welte T100 etwa besaßen 100 Spuren von denen nur 80 zu Tönen zugeordnet waren, die übrigen 20 dienten zur Steuerung des Abspielgerätes \parencite[]{mxp_2002522}.

Solche Informationen konnte im Herstellungsprozess entweder manuell kodiert werden oder aber durch die direkte Aufnahme von Stücken über speziell dafür konstruierte Aufnahmeklaviere erzeugt werden.
An diesen konnten Künstler:innen ein Stück einspielen welches danach in einem mehrstufigen Prozess durch den Rollenhersteller auf eine Notenrolle übertragen wurde.
Das genaue Verfahren mit dem dies Geschah wurde durch die Hersteller mit größter Geheimhaltung geschützt und da keines der Aufnahmegeräte bis heute überdauert hat kann darüber nur spekuliert werden \parencite[]{zoltan_1994}.

Da Audioaufnahmen zur damaligen Zeit noch nicht in vergleichbarer Zeit möglich waren, bescherten grade diese eigens eingespielten Künstlerrollen den Reproduktionsklavieren eine enorme Popularität.
Schnell entwickelte sich eine Vielzahl von verschiedenen Systemen von Reproduktionsklavieren und dazugehörigen Toninformationsträgern und in den 1920er Jahren hatte die Mehrzahl der bekannten Pianist:innen für einen oder mehrere der Hersteller Notenrollen eingespielt \parencite[]{bowers_1972}.

Im Rahmen des vorangegangenen TASTEN-Projekts wurden zwischen 2018 und 2020 am Musikinstrumentenmuseum der Universität Leipzig ca. 5000 Notenrollen in Form von Bildern digitalisiert.
Von diesen liegen ca. 3000 als vollständige Scans vor, von dem restlichen ca. 2000 der Rollen sind nur die Anfänge der Rolle, auf denen Metainformationen wie Titel, Format, etc. vermerkt sind, gescannt worden.
Das gesamte Konvolut umfasst über 30 verschiedene Formate von Notenrollen.
Während im Rahmen des DISKOS-Projektes aus den Scans von Notenrollen aller vorliegenden Formate Midi-Dateien generiert werden sollen, beschränkt sich diese Arbeit auf Rollen des Typs Hupfeld Phonola Solodant, einem verhältnismäßig einfach Format mit 77 Spuren.

% Theorie

\section{Forschungsstand} \label{Forschungsstand}

Das bisherige Forschungsinteresse zur Digtialisierung von historischen Toninformationsträgern konzentriert sich stark auf die Verarbeitung von Notenrollen.
Dies dürfte vor allem auf die Existenz der bereits erwähnten, von kontemporären Pianist:innen eingespielten Notenrollen zurückzuführen sein.
Aber auch die bessere Verfügbarkeit, bedingt durch die hohe produzierte Stückzahl aber auch durch die spätere Produktionszeit und die damit besseren Chancen von heute noch funktionsfähigen Medien, kann ein Faktor sein der zu dieser Diskrepanz am Forschungsinteresse beiträgt.

Die einzigen bekannten Arbeiten zur Digitalisierung von Lochplatten sind im Rahmen des SISAR-Projekts der Italian Mechanical Music Association (AMMI) entstanden \parencite[]{pedrazzini_2013,perretti_2014}.
Im Rahmen dieses Projekts sollen verschieden Typen von historischen Toninformationsträgern digitalisiert werden.
Dabei arbeiten die Autor:innen auch mit Pappplatten des Types Ariston, für die eine Anwendung entwickelt wurde mit der Bilder dieser runden Medien in eine viereckige, einer Notenrolle ähnliche Darstellung transformiert werden können um sie später mit Software zur Digitalisierung von Notenrollen bearbeiten zu können \parencite[]{perretti_2014}.
Eine solche Anwendung wurde ebenfalls im Rahmen des Forschungsprojektes entwickelt \parencite[]{conversion}.
Die Umwandlung in eine Midi-Datei erfolgt dabei mittels hinterlegten Informationen zum Medienformat und den vorhandenen Spuren, das Ausrichten der Spuren auf dem Scan erfolgt manuell.

Leider sind die eingesetzten Methoden praktisch undokumentiert, so dass sich weder zu Qualität und Funktionalität der entstehenden Digitalisate, noch zum aktuellen Entwicklungsstand eine sinnvolle Einschätzung vornehmen lässt.
Auch ist die eingesetzte Software weder als Quellcode noch in binärer Form öffentlich verfügbar, so dass auch auf diesem Weg kein Schluss auf die Güte des Vorgehens möglich ist.

Dennoch können diese Arbeiten zeigen, dass die Digtialisierung von Lochplatten und Notenrollen viele Ähnlichkeiten aufweisen.
Dazu gehören zum einen die ähnliche Art der Informationskodierung (durch Löcher), aber auch viele mechanische Eigenschaften der Abspielgeräte.
Beide Formate wurden durch luftbetriebene Mechanismen von den Abspielgeräten ausgelesen, so dass auch die Art der Informationsdekodierung grundsätzlich ähnlich ist.
Dies legt nah die Digtialisierung der beiden Medienformate als gemeinsamen Themenbereich zu betrachten.
Die vorhandenen Ansätze zur Digitalisierung von Notenrollen können folglich auch als relevant für Verfahren zur Digitalisierung von Lochplatten betrachtet werden.

Die ersten Überlegungen zu solchen Verfahren lassen sich in den 70er Jahren verorten, in den 90ern und frühen 2000ern existierten bereits mehrere Systeme im weiteren Umfeld der digitalen Verarbeitung von Notenrollen \autocite[]{colmenares_2011}.
Allgemein lässt sich feststellen, dass die Projekte aus dieser Zeit überwiegend durch interessierte Privatpersonen und weniger durch akademische Forschung entstanden sind.
Eines der ersten Projekte dieser Art war das \textit{Philips Ampico-Apple} System von Peter Stephens \parencite*[]{stephens}, das die Digitalisierung von Notenrollen des Types Ampico ermöglichte.
Ein späteres Beispiel aus 1996 ist die Entwicklung einer Lösung zur Digitalisierung von Notenrollen mehrerer Formate von Wayne Stahnke \parencite[]{stahnke_1996}.
Als letztes Beispiel dieser privaten Projekte, sei Spencer \textcite[]{chase_2003} genannt, der ebenfalls ein System zur Digitalisierung von Notenrollen konstruierte, aber auch daran interessiert war diese Digitalisate anschließend wieder auf physischen Instrumenten abspielen zu können.
Eine weitere größere Sammlung von digitalisierten Klavierrollen existierte auf der, inzwischen nicht mehr verfügbaren, Website der International Association of Mechanical Music Preservationalists (IAMMP)\footnote{\href{http://www.iammp.org/}{www.iammp.org}}.

Ziel dieser privaten Projekte war es primär die Information auf den Notenrollen zu konservieren und sie auf anderen Geräten wieder abspielbar zu machen.
Gemein haben diese Enthusiasten-Projekte eine sehr beschränkte bis komplett fehlende Dokumentation der Prozesse die zur Digitalisierung genutzt wurden.
Meist wurden die Verfahren entweder entweder gar nicht oder nur auf Mailinglisten oder (inzwischen häufig nicht mehr verfügbaren) privaten Websites dokumentiert.
Die verwendete Software ist häufig proprietär und steht heute nicht mehr für eine Analyse zur Verfügung.
Auch die für die erstellten Digitalisate verwendeten Dateiformate sind insbesondere bei den älteren Projekten proprietär und nur begrenzt dokumentiert.
Bei größeren Sammlungen, wie etwa der der IAMMP, ist eine Nachvollziehbarkeit der Ersteller:innen und der genutzten Methoden der Digitalisate nur eingeschränkt gegeben.
Diese Faktoren schränken die wissenschaftliche Verwendbarkeit der Digitalisierungsprozesse aber auch der erzeugten Digitalisate stark ein.

Eine frühe wissenschaftliche Arbeit zur Digitalisierung von Notenrollen findet sich bei \textcite[]{zoltan_1994}.
Die Autoren beschreiben die Digtialisierung von Notenrollen des Types Welte-Mignon.
Über einen Rollenscanner werden Bilder der Klavierrollen erzeugt aus welchen nach der softwareseitigen Verarbeitung zwei Midi-Dateien generiert werden.
Eine Midi Datei enthält nur die Informationen zu Tonhöhe und Länge, die andere enthält auch die in der Notenrolle kodierten Informationen zu Dynamik ein.
Die Autoren konnten, bedingt durch die verfügbare Technologie, einige Aspekte der Digtialisierung nur bedingt bzw. durch manuelle Korrekturen durchführen, so etwa das Erkennen von Expressionskurven oder der Umgang mit Defekten an der Notenrolle wie Rissen im Papier.
Sie nutzten bereits Techniken die sich in späteren Arbeiten wiederfinden, so wird etwa der Rollenrande dynamisch über den Verlauf der Rolle verfolgt, um durch Deformationen des Papiers entstandene Ungenauigkeiten auszugleichen.

\textcite[]{debrunner_201300} nutzt für die Digitalisierung von Notenrollen des Typs Welte-Philharmonie sowie Welte-Mignon T100 einen Zeilenscanner und eine eigene Software.
Die verwendete Software ist dabei grunsätzlich in der Lage die Digitalisierung in eine Midi-Datei automatisch durchzuführen, benötigt aber bei einem signifikanten Teil der Rollen manuelle Korrekturen.
Das genaue Digitalisierungsverfahren ist nicht dokumentiert, der Schwerpunkt der Arbeit lieft klar auf der Verarbeitung der auf den Notenrollen kodierten Steuerungsinformationen, welche ausführlich dokumentiert werden.

Eine Arbeit die sich spezifisch mit den Herausforderungen der Erzeugung von MIDI Dateien aus bereits digitalisierten, in Matrizenform vorliegenden Klavierrollen beschäftigt findet sich bei \textcite[]{colmenares_2011}.
Die Autor:innen nutzen dabei Digitalisate die mit der Technologie von \textcite[]{stahnke_1996} erzeugt wurden.
Im Gegensatz zu dessen proprietärer Digitalisierungstechnologie setzen die Autor:innen aber einen Schwerpunkt darauf  ihren Umwandlungsprozess ausführlich zu dokumentieren.

Sie gehen dabei sowohl allgemein auf das Dekodierverfahren für Klavierrollen in Matrizenform ein als auch insbesondere auf die Umsetzung der Kontrollspuren ins Midi Format.
Von interesse ist insbesondere das von den Autoren zur Validierung ihrer Ergebnisse eingesetzte Verfahren.
Sie führen dabei einen Wellenformvergleich von, aus den von ihnen erzeugten Midi-Dateien generierten, Audio Dateien mit den von \textcite[]{stahnke_1996} selbst erzeugten und vertriebenen Audio-CDs der gleichen Notenrollen durch \parencite[70ff]{colmenares_2011}.
Dabei können sie die Genauigkeit ihres Verfahrens demonstrieren, das bis auf kleinere Abweichungen authentische Ergebnisse produziert.
Problematisch scheint allerdings, dass die Genauigkeit dieser Validierungsmethode stark von der Genauigkeit der ursprünglichen, von \textcite[]{stahnke_1996} mittels eines prorietären, nicht öffentlich dokumentierten Prozesses durchgeführten, Digitalisierung abhängig ist.

Das aktuellste und umfangreichste Vorhaben zur Digitalisierung von Notenrollen ist das Player Piano Projekt an der Stanford University.
Dieses wurde 2014 ins Leben gerufen um die umfangreiche Sammlung von Notenrollen in der eigenen Bibliothek Forscher:innen und anderen interessierten zugänglich zu machen \autocite[]{shi_2019}.
Aktuell befinden sich circa 20.000 Notenrollen in der Sammlung in Stanford von denen ca. 1000 erfolgreich digitalisiert wurden \autocite[]{broadwell_2022}.
Das Projekt fokussierte sich dabei zunächst auf die Digitalisierung von Notenrollen des Typs Welte-Mignon T-100, umfasst inzwischen aber auch die Digitalisierung Rollen anderer Formate.

Dafür wurde zunächst ein Gerät konstruiert, dass es erlaubt hochauflösende\footnote{TIFF mit 300dpi Auflösung.} Bilder der Notenrollen anzufertigen.
Teil des Scanprozesses ist dabei das anwenden verschiedener, automatischer Fehlerkorrekturen für physische Defekte der gescannten Notenrolle.
So wird etwa, ähnlich zu \textcite[]{zoltan_1994} ein Korrekturverfahren zur Begradigung der Notenrolle genutzt, mit dem, besonders in den späteteren Teilen der Rollen auftretende, Verzerrungen kompensiert werden können.
Weiterhin setzten die Autor:innen Filter ein um Löcher die uncharakteristisch Merkmale aufweisen\footnote{Beispielsweise bei Form oder Orientierung des Lochs.} und vermutlich auf Beschädigungen zurückzuführen sind zu erkennen.
Schnell aufeinanderfolgende Löcher, die von Originalabspielgeräten als einzelne Note gespielt werden\footnote{Sogenannte \textit{Bridges}.}, werden zusammengefasst \autocite[519-520]{shi_2019}.

Auch die Funktionalität der Steuerungsspuren wird in dem von \textcite[521-522]{shi_2019} vorgestellten Digitalisierungsverfahren möglichst originalgetreu in den erzeugten Midi-Dateien abgebildet.
Ebenso werden weitere Anpassungen vorgenommen um das Abspielverhalten auf Originalabspielgeräten zu simulieren, wie etwa die leichte Temposteigerung die durch die sich verändernden Rollendurchmesser beim Abspielen einer Notenrolle vermutet wird.
Aus den generierten Midi-Dateien wird mittels eines generischen Software-Pianos eine Audiodatei generiert.no gerendert wurden.

Aktuell arbeiten die Autor:innen unter anderem am Pianolatron, einer webbasierten Softwarelösung die Notenrollen darstellen und interaktiv abspielen kann \parencite[]{vijoy_2022}.
Die Software erlaubt es dabei eine Notenrolle aus dem Bestand von ca. 1000 digitalisierten Rollen abzuspielen, Parameter wie Abspielgeschwindigkeit und Lautstärke können interaktiv angepasst werden.
Auch können verschiedene Parameter der Emulation der mechanischen Komponenten der Originalabspielgeräte angepasst werden, so dass sich, in Fällen in denen durch Verlust der Originalabspielgeräte keine notwendigerweise originalgetreue Abbildung möglich ist, diese sich experimentell Testen lassen.
Es bestehen weitere Interaktionsmöglichkeiten, so können Nutzer:innen etwa durch Auswählen von auf den Rollen vorhandenen und durch die Software erkannten Noten weiter Informationen zu diesen einsehen.
Insgesamt bildet das Pianolatron schon jetzt eine umfangreiche und vielversprechende Abspiel- und Visualsierungsumgebung für Notenrollen.



% Vorgehen

\section{Vorgehen}



% Umsetzung

\section{Umsetzung}


% Ergebnisse

\section{Diskussion}


% Fazit

\section{Fazit}

Historische Toninformationsträger in Form von u.a. Lochplatten und Notenrollen stellen wertvolle, zeithistorische Zeugnisse aus einer Zeit in der privater Musikkonsum erstmals möglich wurde dar.
Im Rahmen der Forschung am Musikinstrumentenmuseum der Universität Leipzig sollen diese Medien präserviert und für die musikwissenschaftliche Forschung geöffnet werden.
Dazu gehört die softwarebasierte Erzeugung von Midi-Dateien aus Bildaufnahmen der Originalmedien.
Auch sollen verschiedene Analyse- und Visualsierungtools geschaffen werden die zur Beantwortung musikwissenschaftlicher Fragestellungen an die Medien beitragen können.

In dieser Arbeit wurde die Implementierung einer Anwendung zu diesem Zweck in Form der Python Software \code{musicbox} vorgestellt.
Es wurde ein Verfahren eingeführt mit dem diese Anwendung aus Lochplatten erfolgreich und vollständig digitale Midi-Dateien erzeugen kann.
Weiterhin wurde ein grundlegendes Verfahren zur Digitalisierung von Notenrollen vorgestellt, welches in seinen Grundfunktionen bereits für die Digitalisierung einzelner Medien nutzbar ist.
Abschließend wurden erste, in der Anwendung integrierte, Analyse- und Visualsierungstools beschrieben, die zur Beantwortung musikwissenschaftlicher Fragestellungen an den Medienbestand genutzt werden können.

Insgesamt konnte ein erfolgreiches Verfahren für die Digitalisierung historischer Toninformationsträger gezeigt werden.
Die dabei entstandene Softwareanwendung bietet durch ihre auf einfache Erweiterbarkeit ausgelegte Architektur eine solide Grundlage für die Verarbeitung weiterer Medienformate sowie die Umsetzung zusätzlicher Visualisierungs- und Analyseverfahren.
In der Zukunft ist die Erweiterung der Anwendung für diese Zwecke geplant.

\clearpage

%
% ---- Bibliography ----
%
%
\setcounter{biburllcpenalty}{7000}
\setcounter{biburlucpenalty}{8000}

\selectlanguage{german}

\printbibliography[title=Literatur, keyword=literature]
\printbibliography[title=Software, keyword=software]
\printbibliography[title=Online, keyword=online]

\clearpage

% Die eidesstattliche Erklärung mit Unterschrift
\textbf{Erklärung der Urheberschaft}

\vspace{1cm}

Hiermit erkläre ich, daß ich die vorliegende Arbeit
selbstständig angefertigt, nicht anderweitig zu Prüfungszwecken vorgelegt und
keine anderen als die angegebenen Hilfsmittel verwendet habe. Sämtliche 
wissentlich verwendete Textausschnitte, Zitate oder Inhalte anderer Verfasser 
wurden ausdrücklich als solche gekennzeichnet.

\vfill

\hspace{2cm} \textbf{Leipzig, 13.08.2021} \hspace{2cm} \includegraphics[scale=0.3]{graphics/signature.png}


\vspace{1cm}

\hspace{2cm} Ort, Datum \hfill Unterschrift \hspace{2cm}


\end{document}
